%! AP Statistics — Unit 2, Lesson 2.2 — Homework KEY
\documentclass[10pt]{article}
\usepackage{apstats-article}
\usepackage{apstats-key}

%=============================================================================
\begin{document}

\pageheader{Unit 2, Lesson 2.2}{Homework --- Answer Key}
\begin{tabularx}{\linewidth}{X X X}
Name:\;\ans{Key} & Date:\; & Period:\;
\end{tabularx}

\vspace{0.1in}

\begin{tcolorbox}[colback=sky, colframe=navy, boxrule=0.8pt, arc=2mm,
    left=3mm, right=3mm, top=1.5mm, bottom=1.5mm]
\small For each problem: (a) complete missing values, (b) compute conditional
relative frequencies, and (c) state whether there is evidence of an association.
\end{tcolorbox}

\vspace{0.14in}

% ════════════════════════════════════════════════════════════════════════════
% PROBLEM 1 KEY
% ════════════════════════════════════════════════════════════════════════════
\begin{blurbbox}[Problem 1]{navylight}
\small

\vspace{8pt}
\renewcommand{\arraystretch}{1.8}
\begin{tabularx}{\linewidth}{@{} l C{2.4cm} C{2.4cm} C{2.2cm} @{}}
\rowcolor{sky}
                     & \textbf{Remote: Yes} & \textbf{Remote: No} & \textbf{Row Total} \\
\hline
Nexus       & 60  & 40  & 100 \\
Apex        & 30  & 70  & 100 \\
\hline
Col.\ Total & \ans{90} & \ans{110} & \ans{200} \\
\end{tabularx}

\vspace{10pt}
\textbf{(b)} P(Remote $|$ Nexus) $= \ans{60}/\ans{100} = \ans{60\%}$
\hfill
P(Remote $|$ Apex) $= \ans{30}/\ans{100} = \ans{30\%}$

\vspace{8pt}
\textbf{(c)} \ans{Yes, there is evidence of an association. Nexus employees work remotely
at twice the rate of Apex employees (60\% vs.\ 30\%). The conditional distributions differ
substantially, suggesting that company and remote work status are associated.}

\end{blurbbox}

\vspace{0.14in}

% ════════════════════════════════════════════════════════════════════════════
% PROBLEM 2 KEY
% ════════════════════════════════════════════════════════════════════════════
\begin{blurbbox}[Problem 2]{greenacc}
\small

\vspace{8pt}
\renewcommand{\arraystretch}{1.8}
\begin{tabularx}{\linewidth}{@{} l C{2.0cm} C{2.0cm} C{2.0cm} C{2.0cm} @{}}
\rowcolor{sky}
                      & \textbf{Low} & \textbf{Moderate} & \textbf{High} & \textbf{Row Total} \\
\hline
Short sleep  & 18 & 42 & 90 & \ans{150} \\
Adequate     & 60 & 54 & 36 & \ans{150} \\
\hline
Col.\ Total  & \ans{78} & \ans{96} & \ans{126} & \ans{300} \\
\end{tabularx}

\vspace{10pt}
\textbf{(b)} Conditional distributions:

\vspace{6pt}
\renewcommand{\arraystretch}{1.8}
\begin{tabularx}{\linewidth}{@{} l C{2.6cm} C{2.6cm} C{2.6cm} C{1.6cm} @{}}
\rowcolor{sky}
                     & \textbf{Low} & \textbf{Moderate} & \textbf{High} & \textbf{Total} \\
\hline
Short sleep & \ans{12\%} & \ans{28\%} & \ans{60\%} & 100\% \\
Adequate    & \ans{40\%} & \ans{36\%} & \ans{24\%} & 100\% \\
\end{tabularx}

\vspace{10pt}
\textbf{(c)} Segmented bar chart: two bars (Short / Adequate), each reaching 100\%,
divided into Low (bottom), Moderate (middle), High (top). Short sleep: 12/28/60;
Adequate: 40/36/24. High-stress segment is much larger for Short sleep.

\vspace{8pt}
\textbf{(d)} \ans{There is strong evidence of an association between sleep duration and
academic stress. Among short-sleepers, 60\% report high stress, compared to only 24\%
of adequate sleepers. The conditional distributions differ greatly across the two groups.}

\end{blurbbox}

\vspace{0.14in}

% ════════════════════════════════════════════════════════════════════════════
% PROBLEM 3 KEY
% ════════════════════════════════════════════════════════════════════════════
\begin{blurbbox}[Problem 3]{redacc}
\small

\vspace{8pt}
\renewcommand{\arraystretch}{1.8}
\begin{tabularx}{\linewidth}{@{} l C{2.6cm} C{2.6cm} C{2.4cm} @{}}
\rowcolor{sky}
                     & \textbf{Music: Yes} & \textbf{Music: No} & \textbf{Row Total} \\
\hline
Print books & 0.24 & 0.16 & \ans{0.40} \\
E-books     & 0.36 & 0.24 & \ans{0.60} \\
\hline
Col.\ Total & \ans{0.60} & \ans{0.40} & \ans{1.00} \\
\end{tabularx}

\vspace{10pt}
\textbf{(b)} P(Music $|$ Print) $= \ans{0.24}/\ans{0.40} = \ans{0.60 = 60\%}$
\hfill
P(Music $|$ E-books) $= \ans{0.36}/\ans{0.60} = \ans{0.60 = 60\%}$

\vspace{8pt}
\textbf{(c)} \ans{\textbf{No} evidence of association. Both print and e-book readers
listen to music while studying at exactly the same rate (60\%). The conditional
distributions are identical, so there is no association between book preference and
music listening.}

\end{blurbbox}

\vspace{0.14in}

\begin{reflectionbox}
\small
\ans{We use conditional relative frequencies because raw counts and joint relative
frequencies are affected by the sizes of the groups. Conditional frequencies put each
group on the same ``100\%'' scale, allowing fair comparison. For example, in Problem 2
the short-sleep and adequate-sleep groups had the same number of students (150 each),
but the raw count comparison would still be misleading for groups of different sizes.
Joint relative frequencies tell us what fraction of \emph{everyone} falls in a cell,
not what fraction of each group shows the outcome --- that distinction is what reveals
association.}
\end{reflectionbox}

\begin{teachernote}
\textbf{Problem 3 (no association):} A deliberate counterexample --- students should not assume every table shows association. The identical conditional frequencies (60\%/60\%) are the key. Emphasize: association requires the conditional distributions to \emph{differ}.
\textbf{Problem 2 (sleep/stress):} The unequal group sizes (150/150 here, but could differ) reinforce why conditional, not joint, frequencies matter.
\end{teachernote}

\end{document}
